% Generated draft: requires mathematical and source review.
\documentclass[11pt]{article}
\usepackage[margin=25mm]{geometry}
\usepackage{fontspec}
\setmainfont{Latin Modern Roman}
\usepackage{amsmath,amssymb,mathtools,unicode-math}
\setmathfont{Latin Modern Math}
\usepackage{enumitem,needspace,xurl,hyperref}
\setlength{\parindent}{0pt}
\setlength{\parskip}{5pt}
\setlength{\emergencystretch}{3em}
\newcommand{\N}{\mathbb{N}}
\newcommand{\Z}{\mathbb{Z}}
\newcommand{\Q}{\mathbb{Q}}
\newcommand{\R}{\mathbb{R}}
\newcommand{\C}{\mathbb{C}}
\newcommand{\F}{\mathbb{F}}
\newcommand{\A}{\mathbb{A}}
\newcommand{\PP}{\mathbb{P}}
\newcommand{\E}{\mathbb{E}}
\newcommand{\eps}{\varepsilon}
\newcommand{\dd}{\,\mathrm{d}}
\begin{document}
\section*{UnsolvedMath ID: 801152}
\textbf{Source problem code:} catalogue entry
\subsection{Definitions and mathematical statement}
For a conventional smooth formulation, let \(M\) be a smooth manifold of dimension \(n\geq 1\), and let \(L\colon TM\to TM\) be a smooth field of endomorphisms. Its Nijenhuis torsion is
\[
N_L(X,Y)=[LX,LY]-L[LX,Y]-L[X,LY]+L^2[X,Y]
\]
for smooth vector fields \(X,Y\). The operator \(L\) is called Nijenhuis if \(N_L=0\).

On an open set \(U\subseteq M\), a pseudo-Riemannian metric \(g\) is a smooth, symmetric, nondegenerate covariant tensor of order two. Write \(\nabla^g\) for its Levi-Civita connection and \(X^\flat=g(X,\cdot)\). The usual tensor-equation interpretation of geodesic compatibility requires
\[
g(LX,Y)=g(X,LY)
\]
and
\[
\nabla^g_X L
=\frac12\left(X\otimes \dd(\operatorname{tr}L)
+\operatorname{grad}_g(\operatorname{tr}L)\otimes X^\flat\right)
\qquad\text{for every smooth vector field }X\text{ on }U.
\]
Here \((X\otimes\alpha)(Y)=\alpha(Y)X\), and \(\operatorname{grad}_g f\) is characterized by \(g(\operatorname{grad}_g f,Y)=\dd f(Y)\). Let \(\mathcal G_L(U)\) denote the metrics satisfying these equations.

\textbf{Problem.} Characterize the Nijenhuis operators \(L\) that admit at least one geodesically compatible metric, and, for each prescribed \(L\), describe all such metrics. Under the local interpretation, this asks for necessary and sufficient conditions on the germ of \(L\) at an arbitrary point \(p\in M\) for the existence of a neighborhood \(U\ni p\) with \(\mathcal G_L(U)\neq\varnothing\), and for a description of all compatible metric germs at \(p\). The source particularly emphasizes neighborhoods of singular points of \(L\).

\textbf{Formulation qualifications.} The supplied excerpt does not define geodesic compatibility, prescribe a regularity category, specify whether \(L\) must be invertible, or settle whether a global classification is also requested. The smooth tensor-equation formulation above is therefore an explicit conventional interpretation requiring confirmation against the cited source. It imposes no invertibility condition on \(L\). The excerpt also does not define its terms ``singular points'' and ``regular points.''

\textbf{Supplied scope and prior-work qualifications.} The source reports that the existence question is known near almost every point, by the local descriptions in references \texttt{[splitting2]} and \texttt{[splitting1]}. No measure or precise genericity convention for ``almost every'' is supplied. It reports a solution under the additional assumptions that \(L\) has \(n\) distinct eigenvalues at generic points and \(g\) is Riemannian, citing \texttt{[Matveev2006]}, and separately states that the Riemannian case is solved in \texttt{[BMF]}. It describes the case of non-Riemannian signatures as still open even in dimension two, while also citing \texttt{[BMP]} and \texttt{[japan]} for the description of two-dimensional compatible pairs of signature \((+,-)\) near regular points. These are supplied source assessments, not independently verified results; the precise remaining scope, particularly at singular points, requires review.
\subsection{Short English statement}
Determine which Nijenhuis operators admit a geodesically compatible metric and describe every compatible metric for a given operator. The source emphasizes singular points and unresolved questions for indefinite metrics, including dimension two, while reporting known Riemannian and regular-point results. The displayed compatibility equations give a conventional smooth interpretation; the exact definitions, regularity, invertibility requirements, and local or global scope need confirmation from the source.
\subsection{Sources}
\begin{itemize}
\item Local UnsolvedMath registry, permanent ID \(801152\), source problem code \texttt{catalogue entry}, entitled ``Nijenhuis Geometry and Integrable Systems --- Problem 5.'' The supplied registry label is \texttt{open}; this label is not independent verification of mathematical status.
\item Bolsinov, Konyaev, and Matveev, \emph{Research Problems on Nijenhuis Geometry and Integrable Systems} (2024), Problem 5. Supplied source URL: \url{https://arxiv.org/abs/2410.04276}. The record attributes the problem to these authors and gives proposed year 2024.
\item Supplied provenance: extraction method \texttt{arxiv\_latex}; source access date 2026-10-16. This is the access date recorded in the supplied data, not a claim that this draft independently accessed or verified the source. The registry records creation and update on 2026-10-16 and upstream revision \texttt{d371ee917ac9647fbb06cd6fc529deec6f3dda82}.
\item References \texttt{[splitting2]} and \texttt{[splitting1]} are cited in the supplied excerpt for local descriptions of geodesically compatible pairs and the reported answer near almost every point. Their bibliographic entries were not supplied.
\item Reference \texttt{[Matveev2006]} is cited for the reported solution with generically distinct eigenvalues and a Riemannian metric; \texttt{[BMF]} is cited for the reported solution of the Riemannian case. Their bibliographic entries were not supplied.
\item References \texttt{[BMP]} and \texttt{[japan]} are cited for two-dimensional compatible pairs of signature \((+,-)\) near regular points. Their bibliographic entries were not supplied.
\end{itemize}
\end{document}
